\section{Container Architecture (C2)}
The Trade2 backend architecture is organized into containers that work together.
The architecture follows a three-tier approach with clear separation of concerns: presentation (Mobile App), application logic (Backend API), and data storage (Database).
This design supports scalability and maintainability while enabling the integration of monitoring and logging infrastructure required for the \emph{Robustness \& Availability} advanced track.

As this document focuses exclusively on backend architecture, only the Backend API and Database containers are discussed in detail.
The mobile application is only shown for context and is outside the scope of this documentation.

\subsection{Overview}
The mobile application is used by end users to interact with the system.
It is responsible for presenting data and sending requests to the backend.
All communication with the backend happens over HTTP for standard operations and WebSocket connections for messaging.

The backend API is the core of the system.
It handles incoming requests, applies business logic, and coordinates data between the different parts of the system.
It also manages authentication and ensures that users can only access data they are allowed to.

The database is used to persist application data such as users, listings, notifications, and messages.
It ensures that data remains consistent and available between requests.
Object storage from MinIO is used for storing files such as profile images and listing images.

To support the \emph{Robustness \& Availability} track, the backend integrates with a monitoring and logging stack.
Metrics are collected and exposed for scraping by Prometheus.
Logs are aggregated using Loki.
Grafana is used to visualize both logs and metrics, allowing administrators to monitor system health and behaviour.

\Cref{fig:c2-ts} and \cref{fig:c2-ms} show the container diagrams, illustrating how the mobile application, backend, database, storage, and monitoring system interact.

\begin{figure}[H]
    \includegraphics[width=\textwidth]{bd/technical/c2-ts}
    \caption{C2 container diagram (Trade2 System)}
    \label{fig:c2-ts}
\end{figure}

\begin{figure}[H]
    \includegraphics[width=\textwidth]{bd/technical/c2-ms}
    \caption{C2 container diagram (Monitoring System)}
    \label{fig:c2-ms}
\end{figure}

\subsection{Containers}

\subsubsection{Backend API (NestJS)}

\paragraph{Justification}
NestJS was selected based on prototype evaluation comparing NestJS, Spring Boot (Java), and Echo (Go).
Selection criteria included community, developer experience, ecosystem, and quality assurance.

\paragraph{Responsibilities}
\begin{itemize}
    \item Handles incoming HTTP and WebSocket requests.
    \item Implements business logic and enforces access control.
    \item Manages user authentication and authorization.
    \item Coordinates data between Database and Object Storage.
    \item Pushes structured logs to Loki and exposes metrics for Prometheus.
\end{itemize}

\paragraph{Communication}
\begin{itemize}
    \item Receives HTTP and WebSocket requests from the Mobile Application.
    \item Reads and writes data to the Database using PostgreSQL.
    \item Stores and retrieves files from Object Storage (MinIO).
    \item Sends metrics to Prometheus and logs to Loki.
\end{itemize}

\subsubsection{Database (PostgreSQL)}

\paragraph{Justification}
MySQL was considered for familiarity with team skills, but PostgreSQL was chosen for its richer feature set, better support for advanced queries and transactional operations, and extensibility.

\paragraph{Responsibilities}
\begin{itemize}
    \item Persists users, listings, notifications, and messages.
    \item Maintains data consistency and transactional integrity.
    \item Supports queries for search, filtering, and message history.
\end{itemize}

\paragraph{Data Organization}
\begin{htable}[
    caption = {Data Organization},
    label = {tab:data-org}
] {
    colspec = {lX},
}
    Entity & Description \\
    users & User accounts and authentication. \\
    listings & Skill offers and requests with metadata. \\
    notifications & User notifications. \\
    messages & Real-time chat between users. \\
\end{htable}

\subsubsection{Object Storage (MinIO)}

\paragraph{Justification}
MinIO was selected to handle object storage for the Trade2 system due to its compatibility with the S3 API, high performance, scalability, and ease of deployment.
Ceph would have also been a good option, but MinIO was chosen because of team familiarity.
Amazon S3 was considered too, but not chosen due to it introducing a dependency on a cloud provider and its costs.

\paragraph{Responsibilities}
\begin{itemize}
    \item Stores files such as profile images and listing images.
    \item Provides reliable, scalable object storage accessible by the Backend API.
\end{itemize}

\subsubsection{Observability Stack (Prometheus, Loki \& Grafana)}

\paragraph{Responsibilities}
\begin{itemize}
    \item \textbf{Prometheus}: Scrapes metrics from the Backend API and stores time-series data
    \item \textbf{Loki}: Aggregates and indexes structured logs from the Backend API
    \item \textbf{Grafana}: Visualizes both logs and metrics, enabling administrators to monitor system health and behaviour
\end{itemize}
